% -------------------------------------------------------------------------
% ÜBERHOLT (31.08.2026). Maßgeblich ist die eingereichte main.tex. Zahlen,
% Zeilennummern und Angaben zum Overleaf-Stand in dieser Datei stammen aus
% der Zeit VOR der Crosswalk-Korrektur vom 29.08.2026 (35 statt 36 Stadt-
% teile, vollständiger Neulauf) und wurden bewusst NICHT nachgezogen: die
% Datei ist als datierter Entwurf Teil der Arbeitsdokumentation.
% -------------------------------------------------------------------------
% =========================================================================
% KAPITEL 5.2 -- GLEITOBJEKTE
% 24.08.2026. Zu 5.2 gehoert genau EIN Gleitobjekt: Quelltext 1.
% Keine Tabelle, keine Abbildung.
%
% PLATZIERUNG
%   Direkt NACH Absatz 2 (ACS-Join) einfuegen. Der Fliesstext verweist mit
%   \ref{lst:acs-versatz} darauf.
%
% VORAUSSETZUNGEN IM PREAMBLE
%   \lstdefinestyle{pythonstil}{...}   -- laut Schreibplan vorhanden
%   Falls pythonstil anders definiert ist (z. B. \lstdefinelanguage), in der
%   ersten Zeile style= entsprechend anpassen.
%
%   Fuer die deutsche Benennung im Text und im Verzeichnis:
%     \renewcommand{\lstlistingname}{Quelltext}
%     \renewcommand{\lstlistlistingname}{Quellcodeverzeichnis}
%   Der Fliesstext schreibt "Quelltext~\ref{lst:acs-versatz}". Ohne das
%   \renewcommand heisst es im Verzeichnis "Listing" -- dann entweder das
%   \renewcommand setzen oder im Text "Listing" schreiben. Nicht mischen.
%
%   Die literate-Zuordnung fuer Umlaute wird hier NICHT gebraucht: Der
%   Ausschnitt ist reines ASCII. Ab Quelltext 6 (tests/) ist sie noetig.
%
% BESCHRIFTUNG
%   Bewusst einzeilig. Der Satz "die Auswahl greift auf den juengsten
%   veroeffentlichten Jahrgang zurueck" stand hier doppelt -- er steht schon
%   im Fliesstext von Absatz 2 und ist dort besser aufgehoben. Die
%   Beschriftung traegt nur noch Gegenstand und Fundstelle.
%
% QUELLE, ZEICHENGLEICH
%   prep/s1_daten.py, Z. 306--322 (Funktionskopf und Rueckgabe) und Z. 337
%   (Aufruf). Gekuerzt ausschliesslich durch Weglassen ganzer Zeilen:
%   entfernt ist der Docstring. Keine Zeile ist umformuliert.
% =========================================================================


% ---------------------------- AB HIER KOPIEREN ---------------------------

\begin{lstlisting}[style=pythonstil,
  label={lst:acs-versatz},
  caption={Zeitbewusste Zuordnung des ACS-Jahrgangs
           (\texttt{prep/s1\_daten.py}, Z.~306--322 und 337, gekürzt)}]
def acs_snapshot(jahr: int, acs_years: list[int]) -> int:
    return max([a for a in acs_years if a <= int(jahr) - ACS_PUBLIKATIONS_LAG],
               default=min(acs_years))

sffd["acs_year"] = sffd["year"].apply(lambda y: acs_snapshot(y, jahrgaenge))
\end{lstlisting}

% ---------------------------- BIS HIER KOPIEREN --------------------------
